\section{Unconstrained sampling and scalar transfer}
\label{sampling:section}

We compare averages on the parent box with averages on random affine
boxes.  Proposition~\ref{sampling:scalar-transfer} transfers a scalar
comparison against a fixed niltest.  Its proof first stabilizes the
niltest under conditioning on small-prime residues, then uses contraction
on the remaining residue coordinates.

Section~\ref{sampling:constrained-section} constructs the affine
sampling law when polynomial constraints are imposed.
Sections~\ref{cube:section} and~\ref{cube:transfer-section} then treat
path-dependent tests and positive score recovery. Here the structured
test is fixed, and the purpose is to preserve a scalar comparison.

All path probabilities and expected testing suprema refer to the
ensemble constructed before productive paths are selected.  A subsequent
discard is estimated in that original probability space.  Conditioning
on productivity does not supply a new residue or equidistribution law.

\subsection{Scalar comparison without constraints}

\begin{proposition}[Scalar transfer]
\label{sampling:scalar-transfer}
Fix the degree, an integer $\ell\ge2$, and $\varepsilon>0$.  There are
constants $C,C'>0$ with the following property;
$C'$ depends only on $\ell$, and one may take $C'=\ell+1$.
Let $p\ge2$, let $\mathcal I$ be an integer box of dimension $n\le p$ and
side lengths $H_i$, and let $g:\mathcal I\to[0,1]$ be a niltest of
complexity at most $p$.  Let $0\le f\le1$ and $e^{-p}\le\lambda\le2$.
Suppose, for every parameter slice $B\subseteq\mathcal I$ of log cost
at most $P=(2+p)^C$, that
\begin{equation}\label{sampling:no-increment}
 \E_B fg\le\lambda\E_Bg+e^{-P}.
\end{equation}
Extend $fg$ and $g$ by zero.  Choose integer $x,V$ uniformly from product
boxes at scales $H_i$ and $H_i/L$, respectively, with $V$ having $\ell$
columns.  Sample $t$ uniformly from an integer box in
$[-e^pL,e^pL]^\ell$, possibly restricted to one residue in each coordinate
modulo a common $D\le e^p$, with every restricted side at least $Le^{-2p}$.
If $L\ge\mathcal B(p)$ and $\min_iH_i\ge\mathcal B(p)L^{C'}$, then
\begin{equation}\label{sampling:scalar-conclusion}
 \Pr_{x,V}\left\{
 \E_t\bigl(fg-(1+\varepsilon)\lambda g\bigr)(x+Vt)>e^{-p}
 \right\}\le e^{-p}.
\end{equation}
The same conclusion holds when the outer law of $(x,V)$ is dominated by
these product laws by $e^p$, and for scale changes by $e^{O(p)}$, after polynomial changes
of the log budgets.  All complexity and precision requirements are prescribed
before choosing $L$.  The conditional law of $t$ remains the stated uniform
law on the prescribed parameter slice; arbitrary reweighting of $t$ after
observing $(x,V)$ is not included.
\end{proposition}

\subsubsection{Stability of the structured baseline}

The no-increment hypothesis compares the mass of $fg$ with the mass
of $g$ on each slice.  Conditioning on a few residue coordinates can
change both masses, so the reference mass must be controlled before
we use contraction on the remaining coordinates.  The next proposition
does this for the structured function $g$.  It gives a partition on
which every further restriction involving a prescribed number of
prime coordinates changes the mean of $g$ by at most $\delta$.

\begin{proposition}[Stability of baseline marginals]
\label{sampling:stability}
Let $g$ be as in Proposition~\ref{sampling:scalar-transfer}.  Select a finite
collection of distinct primes $l$, with one power $q_l=l^{a_l}$ per prime,
where $\log q_l$ and the logarithm of the number of primes have prescribed
polynomial bounds.  Regard $u\bmod q_l\in(\Z/q_l\Z)^n$ as a single
coordinate.  Prescribe $\delta>0$, an integer $r_0$, and an initial set of
mandatory coordinates, with $\log\delta^{-1}$, $r_0$, and the number of
mandatory coordinates polynomially bounded.

On any prescribed box of polynomial log cost, there is a decision tree
whose nodes fix entire prime coordinates,
with every leaf fixing all the prescribed mandatory coordinates,
whose depth on every branch is polynomially bounded, and such that every
nonempty leaf $B$ satisfies
\begin{equation}\label{sampling:stable-leaf}
 \left|\E[g\mid B,\,u_J=z]-\E[g\mid B]\right|\le\delta
 \qquad(|J|\le r_0)
\end{equation}
for every set $J$ of still unfixed coordinates and every nonempty refinement
$u_J=z$.  The depth and the sufficient lower bound for the original side
lengths have polynomial log bounds in the prescribed inputs, independently
of the coefficients of the polynomial orbit defining $g$.
\end{proposition}

\begin{proof}
We use Lemmas~\ref{prelim:vertical}, \ref{prelim:step-drop},
\ref{prelim:symbol-splitting}, and \ref{prelim:reconstruction}, together
with the actual-factor lifting and cover constructions in
Section~\ref{prelim:section}.  Let $A\ge2$ dominate the original
complexity, $\log\delta^{-1}$, $r_0$, the number of mandatory primes,
$\max_l\log q_l$, the logarithm of the number of selected primes, and
the initial slice cost.  Replace $\delta$ by $\min(\delta,1)$ if
necessary.  First split on all mandatory coordinates.  If the
starting box is a parameter slice with residue modulus $M_{\rm init}$,
also fix every selected prime dividing $M_{\rm init}$; this adds at
most $\log_2 M_{\rm init}$ coordinates.  The initial residue modulus
is then coprime to every set of still-unfixed prime coordinates.

We first describe how a failure of stability produces new rational
information.  At a positive degree $j$, the comparison concerns a
fixed polynomial map $a$ into a filtered group $G$.  Retain a graded
subalgebra $V$, initially the full associated graded algebra, obtained
by intersecting the fast algebras supplied by earlier failures.
Its role is to record the part of the symbol not yet eliminated.
There are two possibilities: the present discrepancy forces a strict
decrease of $V$, or it passes to a map of smaller degree.  Since a
degree-zero map cannot distinguish the two laws, some positive-degree
subalgebra must eventually decrease.

Here is the reduction.  Suppose the two laws at a node and a residue
refinement differ on the incoming test by $d_*>0$.  Approximate that
test uniformly to error $d_*/100$ by top-vertical terms.  Let $M$
bound both their number and total coefficient mass; $\log M$ is
polynomial in the current comparison budget.  Call a term
significant when its average on either law has modulus at least
$d_*/(100M)$.  If a significant frequency $\eta$ does not annihilate
$V_j$, step drop on that law's residue box gives a fast algebra $U$
with $\eta(U_j)=0$.  Consequently
\[
                         \dim(V\cap U)<\dim V.
\]
Record this splitting and replace $V$ by $V\cap U$.  Fix the witness
coordinates $J$ to their full selected prime powers, splitting into
all assignments rather than retaining only the assignment that
witnessed the discrepancy.  Translation does not change the symbol,
so the same record applies on every sibling after transport to its
coordinates.  We call this strict decrease an event.

Otherwise every significant frequency annihilates $V_j$.  Averaging
over the central torus generated by this top subspace loses at most $d_*/100$
on each law, apart from the Fourier approximation error.  After
factoring the map and freezing its controlled slow and rational
factors, the surviving discrepancy passes to a map of degree at
most $j-1$.  The constructions below justify this descent with one
fixed lower-degree map.  Repeating the alternatives must reach an
event, since the discrepancy cannot persist at degree zero.

It remains to bound the cost of repeating this procedure.
Refactoring the map produced by each successive violation would
compose polynomial complexity bounds once per violation, potentially
more than a fixed number of times.  We instead retain the map and
combine its recorded symbol restrictions in one simultaneous
splitting, so that a new polynomial is composed into the budget only
when the degree decreases.  A new
violation starts from its own node and refinement, but uses the same
retained maps until a recorded subalgebra changes.  In particular,
the subboxes used to freeze a slow factor are chosen afresh from
that factor; they are not intersections of the subboxes selected
for earlier violations. This controls the accumulated costs; the
strict dimension decreases will bound the number of events. We now
record the quantitative invariant and verify the factorization,
transport, and descent used above.

\paragraph{The fixed-map invariant.}
Maintain at most $s$ levels, of strictly decreasing filtered degree.
An epoch at degree $j$ has a fixed polynomial map $a$ into a fixed
filtered group $G$, expressed in the affine integer coordinates of
a fixed ancestor residue box.  Its coefficients are unrestricted.
The physical box underlying this ancestor has sides $W_i$.
At a node, a comparison at this level evaluates this same map under
the node and a refinement by still-unfixed primes $J$, possibly on
an additional subbox and one auxiliary residue class.  The subbox
has physical sides at least $e^{-B_j}W_i$, the auxiliary modulus
has logarithm at most $B_j$ and is coprime to the primes in $J$,
and the test complexity and inverse discrepancy logarithm are at
most $B_j$.  The bounds depend on the parent epoch, not on the
number of visits to this level.

At this level retain the intersection $V$ of the recorded graded
fast subalgebras.  For each new value of $V$, retain one factorization
and one child map.  An event at this level resets the lower levels;
an event above it resets this entire epoch.  Each new violation of
\eqref{sampling:stable-leaf} starts again from its node/refinement
pair and descends through the retained maps.  Auxiliary subboxes
from a previous violation are not successively intersected.

\paragraph{Transport and joint intersection.}
Write the old piece and current node in physical coordinates as
\[
 u=a_0+M_0d_0t,\qquad u=a_1+M_1v,\qquad M_0\mid M_1,
\]
where $M_0,M_1$ are the accumulated residue moduli and
$d_0$ is the old controlled auxiliary modulus.  Thus
$t=c+hv$, with $h=(M_1/M_0)/d_0$.
The homogeneous-polynomial identity for symbols removes $c$ and
multiplies a coefficient of multidegree $\alpha$ by $h^{|\alpha|}$;
the identity holds for this rational $h$ as well as integer dilations.
If an old rational symbol has common denominator $D_0$, its
transported common denominator divides $D_0d_0^s$.
The integer $M_1/M_0$ affects only numerators, and the fast
subalgebra's rational orientation is unchanged.

Let $T_i,T'_i$ be old and current parameter sides.  If the old
physical piece has sides at least $\rho W_i$ and the current
physical span is contained in the epoch box, then
\[
 |c_\alpha h^{|\alpha|}|(T')^\alpha
 \le \mathcal B(B_j)
 \prod_i\left(\frac{M_1T'_i}{M_0d_0T_i}\right)^{\alpha_i}
 \le \mathcal B(B_j)\rho^{-s}.
\]
Thus every recorded slow symbol is controlled on the current full
residue box, with no accumulated-modulus or aspect-ratio loss.
There are at most $\dim G$ records in this epoch.  Multiply their
common denominators if needed and apply the simultaneous assertion
of Lemma~\ref{prelim:symbol-splitting} jointly.  The resulting common
splitting has polynomial logarithmic cost.  No repeated powering
of a complexity bound at the same level is used.

\paragraph{The actual factorization and fixed child.}
After each event, on each resulting full residue box write the
current map as $a$ in its affine coordinates and let
$X=E_{\rm sym}PR_{\rm sym}$ be its common symbol splitting, with
$P$ taking values in $V$.  Lift the homogeneous logarithmic
coefficients of $E_{\rm sym},R_{\rm sym}$ through fixed rational
sections, adding no other logarithmic terms; call the lifts
$E_{\rm hom},R_{\rm hom}$.  Choose $\gamma\in\Gamma$ such that
$c=a(0)\gamma$ is a bounded representative, and define
\[
 E=cE_{\rm hom},\qquad R=R_{\rm hom}\gamma^{-1},
 \qquad b=E^{-1}aR^{-1}.
\]
Then $a=EbR$, $b(0)=1$, and the symbol of $b$ lies in $V$.
The finite BCH formula shows that $E$ is fully slow: its constant
is bounded, and every product of nonconstant slow coefficients
has the corresponding product of inverse-side factors.
It also shows that $R$ has one controlled common denominator,
since $\log\Gamma\subset l^{-1}\mathbb Z^d$, even though the
numerators of $\gamma$ may be arbitrarily large.

Set
\[
 \mathfrak h_i=\{v\in\mathfrak g_i:
             v\bmod\mathfrak g_{i+1}\in V_i\}.
\]
Graded bracket closure makes these a filtration.  The degree-$i$
coefficients of $\log b$ lie in $\mathfrak h_i$ by its symbol
condition, while coefficients of larger degree already lie in
$\mathfrak g_{i+1}\subset\mathfrak h_i$.
Hence $b$ takes values in $H=\exp\mathfrak h_1$ and is adapted to
this filtration.  The map $b\bmod H_j$ has degree at most $j-1$.
No bound on the lower-degree coefficients of $b$ is asserted or used.

Let $D_R$ be the actual common denominator of this fixed polynomial
$R$, enlarged by the fixed structural and lattice denominators.
At integer arguments, its values and their adjoint matrices have
denominators dividing a fixed controlled power of $D_R$.
The cover construction in Section~\ref{prelim:section} therefore
gives one controlled lattice $\Lambda$ contained in every required
$R(t)\Gamma R(t)^{-1}$.  Its cost depends on $D_R$, not on the
numerators or on $t$.  We use the actual common denominator here;
we do not take a common multiple of every integer below its bound.
The source lattice $H\cap\Lambda$ and the quotient data for
$H\to H/H_j$ are now fixed.

The same construction gives a controlled period $D$ for $R$ in
right lattice class.  Include every selected prime dividing $D$,
to its full selected power, keeping the child map and cover
unchanged while splitting these assignments.  At most $\log_2D$
coordinates are added.  The selected power may be smaller than
the corresponding prime power in $D$, so this need not make $R$
constant.  It does ensure that every remaining auxiliary period
is coprime to every future $J$: its selected primes have been
removed from the unfixed collection.  Period primes absent from
the collection can never belong to $J$.
When $V$ is full, the fixed child is simply the ordinary top quotient.

\paragraph{Descent when there is no event.}
Average the actual incoming test over the central torus associated
with $H_j$, obtaining $\bar F$.  This removes exactly the vertical
frequencies not annihilating $V_j$.  If none of them is significant,
their total average on each law is at most $d_*/100$.
Including the Fourier approximation errors, $\bar F$ retains at
least a fixed fraction of $d_*$ as discrepancy.  It remains
nonnegative, bounded and Lipschitz.

Use the retained factorization to freeze $E$ on small subboxes
and $R$ exactly in right lattice class on residue pieces.  The
freezing error is independent of $b$, by right invariance:
\[
 d_G(E(t)b(t)R(t),E_0b(t)R(t))=d_G(E(t),E_0).
\]
A fully slow polynomial remains fully slow on later affine
restrictions within its creation box: expansion of a translated
monomial uses bounded normalized offsets, and dilation cancels
against the new sides.  Thus subboxes of controlled relative
physical sides suffice throughout this epoch.  They may be chosen
with comparable lengths; rounding or boundary remainders have
arbitrarily small aggregate mass once the sides are large enough.
At each visit these pieces are obtained afresh from the retained
factors, which proves the fixed relative-side part of the invariant.

For a frozen value $R_0$, the function
\[
 h\longmapsto\bar F(E_0hR_0\Gamma)
\]
is well defined on $H/(H\cap\Lambda)$.  Fixed right translation
is an isometry, bounded left translation by $E_0$ has controlled
Lipschitz norm, and the rational subgroup metric comparison is
controlled.  Since $H_j\subset G_j$ is central, this function is
invariant under left $H_j$.  Lemma~\ref{prelim:reconstruction}
applies to the rational quotient $H\to H/H_j$, whose kernel is
connected and rational and whose differential is controlled.
It gives a controlled nonnegative test of the already fixed
child map.  The target cover can be chosen from the fixed source
lattice and quotient map, independently of the frozen values.

It remains to compare the weights of the pieces.  If the current
selected-prime modulus is $M$, the pending refinement modulus is
$m_J$, and the auxiliary piece modulus is $d_0$, then
$\gcd(m_J,Md_0)=1$.  The $M$- and $d_0$-residues may overlap, but
their compatibility is unchanged by the $J$-refinement.
Counting integers in each physical side modulo the appropriate
least common multiple gives its length divided by that modulus,
with error at most one.  Sufficiently large sides therefore make
the two normalized piece laws arbitrarily close in total variation,
uniformly in their residue assignments.  After the freezing and
weight errors, some piece has normalized discrepancy at least
$d_*/2$, after harmless adjustment of the precision constants.
There is no loss by the number of pieces: the discrepancy is a
weighted average of their discrepancies.

The new test, auxiliary modulus, inverse relative sides and
inverse discrepancy logarithm are polynomially bounded in the
current epoch's bounds.  This is the invariant at the next level.
At degree zero the map is constant, so a positive discrepancy is
impossible.  Every original violation therefore forces an event
at some positive degree.

\paragraph{Budgets, depth and termination.}
Choose $\Phi(x)=(2+x)^{C_s}$ to dominate the fixed number of
Fourier, step-drop, joint-intersection, lifting, cover and freezing
costs just used.  Before constructing the tree set
\[
 B_s=\Phi(A),\qquad B_{j-1}=\Phi(B_j).
\]
There are only $s$ compositions.  At degree $j$, all step-drop
inputs depend on the fixed incoming bound $B_j$, not on the
common splitting assembled from that level's previous records.
Only the child budget depends on this common splitting.
Consequently all $B_j$ are polynomial in $A$.

Every child is a subgroup quotient and covers preserve dimension.
If the original dimension is $d$, the number $N_j$ of events
through a degree-$j$ stack satisfies the coarse recurrence
\[
 N_0=0,\qquad N_j\le d+(d+1)N_{j-1},
 \qquad N_s\le(d+1)^s-1.
\]
Indeed a level has at most $d$ strict dimension decreases and
therefore at most $d+1$ successive child epochs.  Each event adds
at most $r_0$ witness coordinates and polynomially many period
coordinates.  Including the mandatory coordinates, branch depth
is polynomial in $A$.  Each selected $\log q_l$ is bounded by $A$,
so the accumulated branch modulus also has polynomial logarithm.

All auxiliary moduli, inverse relative sides and required inverse
errors have polynomial logarithms.  Thus an original side bound
$\min_iH_i\ge\exp(\operatorname{poly}_s(A))$ suffices simultaneously
for every step-drop, separation and residue-counting use, including
all refinements of size at most $r_0$.
There is no circular choice: choose this bound for a construction
truncated after the displayed maximum number of events; the
dimension recurrence rules out the next event.  The resulting
finite tree terminates, and a leaf cannot violate
\eqref{sampling:stable-leaf}, since a violation would force another
event.
\end{proof}

\begin{remark}[Uniformity of the stability estimates]
Lemma~\ref{prelim:step-drop} gives bounds independent of polynomial-orbit
coefficients, Lemma~\ref{prelim:symbol-splitting} performs simultaneous
splitting by joint linear algebra, and Lemma~\ref{prelim:reconstruction}
controls the reconstructed test on rational images.  The fixed-map and
cover constructions in the proof above preserve these bounds for the
frozen right-translate test independently of the accumulated prime modulus.
\end{remark}

\subsubsection{Residue contraction and low-degree projections}

The residue operator contracts a component once for every prime on
which it depends. Components involving many primes are therefore small.
For the components involving only a few primes, we need bounds from
their conditional marginals; stability and the no-increment hypothesis
will supply those bounds. The lemmas below separate these two estimates.

Residues of an integer point are close to product measure on each
small collection of prime coordinates, but not necessarily on the
entire collection. We therefore also prove versions of the low-degree
estimates that use only these small marginals.

\begin{lemma}[A single-prime contraction]\label{sampling:prime-contraction}
Let $q=l^a$, let $X\in(\Z/q\Z)^n$ and
$V\in(\Z/q\Z)^{n\times\ell}$ be independent uniform variables, and let
$t\in(\Z/q\Z)^\ell$ be independent and uniform.  The operator
\[
 (T\varphi)(X,V)=\E_t\varphi(X+Vt)
\]
from functions on $(\Z/q\Z)^n$ to functions of $(X,V)$ has norm at most
$l^{-\ell/2}$ on the subspace of mean-zero functions.
\end{lemma}
\begin{proof}
Expand in the orthonormal characters
$\chi_\xi(u)=e(\xi\cdot u/q)$.  For $\xi\ne0$,
\[
 T\chi_\xi(X,V)=\chi_\xi(X)
 1_{\{\xi\cdot V_j=0\ \text{for every }j\}}.
\]
The image of the homomorphism $v\mapsto\xi\cdot v$ has order at least
$l$.  Thus $\|T\chi_\xi\|_2^2\le l^{-\ell}$.
The $X$ average makes the images of different characters orthogonal, so
Parseval proves the claim.
\end{proof}

The low-degree global-function strategy here is motivated by global
hypercontractivity; see \cite{KeevashLifshitzLongMinzer2021} for background.
The bounded-marginal estimate below and its approximate-product extension
are proved here.

A stronger bound for an exact product law follows from
Keller, Lifshitz and Marcus~\cite[Theorem~5.4]{KellerLifshitzMarcus2026}.
We give a moment proof of the bound needed here, because the same
argument will also handle approximately product laws.

\begin{lemma}[Low-degree mass under bounded marginals]\label{sampling:low-levels}
Let $\Omega=\prod_{i=1}^m\Omega_i$ carry a product probability measure $\mu$,
where the coordinate spaces are finite. Let $F:\Omega\to[0,M]$ and $K\ge1$.
Suppose
\[
  \mathbb E[F\mid x_J]\le K^{|J|}
       \quad (J\subseteq[m],\ |J|\le r),
\]
pointwise on positive-mass atoms, including $\mathbb EF\le1$ when $J=\varnothing$.
Define the orthogonal product components (also called ANOVA components) by
\[
 F_S=\sum_{T\subseteq S}(-1)^{|S|-|T|}\E[F\mid x_T].
\]
Thus $F=\sum_SF_S$, and each nonconstant component is centered in every
coordinate on which it depends. Put
\[
 U_k^2=\sum_{|S|=k}\|F_S\|_2^2.
\]
If $p_1\ge\max\{r,\log(2+M)\}$, then, with
$R=4(1+K)(p_1+2)$,
\[
 U_k\le R^{2k}\qquad(0\le k\le r).
\]
\end{lemma}

\begin{proof}
We first record an elementary mixed-norm fact. Suppose a function $A$ of $k$
variables, on arbitrary measure spaces, has $L^2$ norm at most $E_0$, and every
section obtained by fixing any $j$ variables has $L^2$ norm at most $E_j$.
In particular $\|A\|_\infty\le E_k$. If
$b_1\ge\cdots\ge b_k\ge2$, with variable $1$ outermost in the mixed norm,
then
\begin{equation}\label{sampling:mixed-norm}
 \|A\|_{L^{b_1}\cdots L^{b_k}}
 \le E_0^{2/b_1}
 \prod_{j=1}^{k-1}E_j^{2/b_{j+1}-2/b_j}
 E_k^{1-2/b_k}.
\end{equation}
To see this, interpolate the innermost norm between $L^2$ and $L^\infty$.
For $B=(\int|A|^2\,dx_k)^{1/2}$ the remaining norm is bounded by
$E_k^{1-2/b_k}$ times the mixed norm of $B$, with exponents
$2b_1/b_k,\ldots,2b_{k-1}/b_k$, raised to the power $2/b_k$.
The section bounds for $B$ are $E_0,\ldots,E_{k-1}$, by Fubini.
Induction on $k$ proves \eqref{sampling:mixed-norm}. Zero bounds follow by
approximation or by observing that the relevant section vanishes.

We prove the lemma by induction on $k$, simultaneously over all finite
product spaces and all densities satisfying the hypotheses. For $k=0$,
$U_0=\mathbb EF\le1$. For the induction step, represent the array of
$F_S$, $|S|=k$, by the tensor
\[
 A((i_1,z_1),\ldots,(i_k,z_k))
 =\begin{cases}
 F_{\{i_1,\ldots,i_k\}}(z_1,\ldots,z_k),&i_1<\cdots<i_k,\\
 0,&\text{otherwise},
 \end{cases}
\]
on the disjoint union of the coordinate spaces, with measure
$\nu=\sum_i\mu_i$. Its $L^2$ norm is exactly $U_k$.

Fix any $j>0$ tensor variables. If their coordinate indices are repeated
or incompatible with the ordering, the section vanishes. Otherwise call
their set $T$, and fix their values $z_T$. For $A_0\subseteq T$, define a
density on the coordinates outside $T$ by holding $x_{A_0}=z_{A_0}$ and
averaging $F$ in $x_{T\setminus A_0}$; denote this density by $F^{A_0}$.
The density $F^{A_0}/K^{|A_0|}$ has cap at most $M$ and satisfies the
marginal hypotheses through level $r-|A_0|$. Consequently its
level-$(k-j)$ norm is at most $R^{2(k-j)}$ by induction. Expanding the
centering operators in the fixed variables gives
\[
 F_{T\cup U}(z_T,x_U)
 =\sum_{A_0\subseteq T}(-1)^{j-|A_0|}(F^{A_0})_U(x_U),
 \qquad |U|=k-j,\quad U\cap T=\varnothing.
\]
The ordering restriction merely omits some $U$ from the sum of squared
norms. The triangle inequality in that direct-sum Hilbert space therefore
gives the uniform section bound
\[
 E_j=(1+K)^j R^{2(k-j)}\le R^{2k-j}\le R^{2k-1}
       \qquad(1\le j\le k).
\]
Use $E_0=U_k$.

Put $P_k=\sum_{|S|=k}F_S$, and choose an even integer $q$ with
$p_1\le q\le p_1+2$. In the expansion of $\mathbb E P_k^q$, a term
vanishes if a coordinate index occurs exactly once, because each ANOVA
component is separately centered. Group the remaining terms by the
partition of the $qk$ slots into blocks of equal coordinate indices.
There are at most $(qk)^{qk}$ such partitions, and every block has
multiplicity at least two. A partition repeating a block within one tensor
copy contributes zero, by the strict ordering of that copy's indices;
discard these partitions. For each remaining partition, take absolute values and
drop the requirement that distinct blocks receive distinct indices.
Its contribution is then bounded by an integral of $q$ copies of $|A|$
over independent variables in the disjoint-union measure $\nu$.

Order these block variables by decreasing multiplicity. Successive
H\"older inequalities, starting with the innermost variable, give one
mixed norm of $A$ per tensor copy; in each norm the exponent of a variable
is its block multiplicity. Indeed, at a block of multiplicity $b$, exactly
$b$ tensor copies involve that variable, so H\"older uses exponent $b$
for each of those copies. The exponents in each resulting tensor norm
are nonincreasing from outermost to innermost. By the mixed-norm bound, each norm is at
most $\max_{0\le j\le k}E_j$, since the displayed powers in the mixed-norm bound are
nonnegative and sum to one. It follows that
\[
 \|P_k\|_q\le(qk)^k\max\{U_k,R^{2k-1}\}.
\]

Orthogonality gives $U_k^2=\mathbb E FP_k$. Also
$\|F\|_{q/(q-1)}\le M^{1/q}(\mathbb EF)^{1-1/q}\le e$.
Since $qk\le(p_1+2)^2$ and $k\ge1$, our choice of $R$ implies
$e(qk)^k\le R^{2k}$. Thus the moment bound yields
\[
 U_k^2\le R^{2k}\max\{U_k,R^{2k-1}\}.
\]
If the maximum is $U_k$, divide by $U_k$ when it is nonzero. Otherwise
$U_k\le R^{2k-1/2}$. In either case $U_k\le R^{2k}$, completing the
induction.
\end{proof}


\begin{lemma}[Truncated projections for an approximately product law]
\label{sampling:gram}
Let $\Omega=\prod_{a\in\mathcal A}\Omega_a$ have a product probability law
$\mu$, and let $\nu$ be another probability law on $\Omega$.  Suppose that
for every $S\subseteq\mathcal A$ with $|S|\le2b$, the density
$d\nu_S/d\mu_S$ differs from $1$ by at most $\eta$ pointwise.
Let $0\le F\le M$ and let $h=d(F\nu)/d\mu$.  Let $h_S$ be its orthogonal
product projections, and put $P=\sum_{|S|\le b}h_S$ and
$A=\sum_{|S|\le b}\|h_S\|_{L^2(\mu)}^2$.  If
$N_b=\sum_{j=0}^b\binom{|\mathcal A|}{j}$, then
\[
 A\le 2M\E_\nu F+2\eta N_b^2 4^bM^2(1+\eta)^2.
\]
\end{lemma}
\begin{proof}
The marginal density of $F\nu$ on any $S$ of size at most $b$ is bounded
by $M(1+\eta)$.  Inclusion--exclusion therefore gives
$\|h_S\|_\infty\le2^{|S|}M(1+\eta)$.
Orthogonality for $\mu$ gives $\int hP\,d\mu=A=\int P^2\,d\mu$.
Each term in $P^2$ uses at most $2b$ coordinates, so
\[
 \left|\E_\nu P^2-A\right|
 \le\eta N_b^2 4^bM^2(1+\eta)^2=:E.
\]
Consequently Cauchy--Schwarz and $F^2\le MF$ give
$A^2\le M\E_\nu F\,(A+E)$.  Writing $B=M\E_\nu F$, this implies
$A\le B+\sqrt{BE}\le2B+2E$, as asserted.
The same proof applies to a finite positive base measure, as will be
needed after conditioning by reference atom masses.
\end{proof}

\begin{lemma}[Low levels for an approximately product base measure]
\label{sampling:approx-low-levels}
Let $\mu$ be a product probability measure on $N$ finite spaces,
and let $\nu$ be a finite positive measure on the same product.  Let
$0\le F\le M$, $K\ge1$, and $h=d(F\nu)/d\mu$.  Suppose
\[
 \E_\mu[h\mid x_J]\le K^{|J|}\qquad(|J|\le r),
\]
including $\int h\,d\mu\le1$.  Let $p_1\ge\max\{r,\log(2+M)\}$,
and choose an even integer $q$ with $p_1\le q\le p_1+2$ and $q\ge2$.
Assume every marginal of $\nu$ on at most $r(q+1)$ coordinates has density
within $\eta$ of $1$ relative to $\mu$, where
\[
 \eta\le\tfrac12\bigl(2^{r+1}(2+N)^r(2+M)\bigr)^{-q}.
\]
Then, with $R=8(1+K)(p_1+2)$,
\[
 \left(\sum_{|S|=k}\|h_S\|_{L^2(\mu)}^2\right)^{1/2}
 \le R^{2k}\qquad(0\le k\le r).
\]
The assertion remains uniform after conditioning or marginalizing
coordinates, provided their number plus the needed marginal order is at
most the original cutoff.  Here conditioning a measure on $x_A=z$ means
restriction divided by its \emph{reference} atom mass $\mu_A(z)$; the
resulting base measure need not have mass exactly one.  The weighted
marginal hypotheses are renormalized by $K^{|A|}$ after fixing $A$.
\end{lemma}
\begin{proof}
Repeat the induction in Lemma~\ref{sampling:low-levels} for the product
projections of $h$.  We give the changes explicitly.  Restricting a set
$A$ of coordinates to fixed values and marginalizing any other coordinates
produces a base measure whose low-order densities are the corresponding
marginal densities of $\nu$ divided by the reference atom mass.  They
still lie between $1-\eta$ and $1+\eta$.  Its weight is a conditional
expectation of $F$, hence remains between zero and $M$.  The weighted
marginal hypotheses, divided by $K^{|A|}$, remain exact.  Thus the
section induction has the same bounds
$E_j=(1+K)^jR^{2(k-j)}\le R^{2k-j}$ for $j>0$, and $E_0=U_k$.
Along a branch at most $r$ coordinates are fixed; the assumed marginal
order leaves at least $qk$ free coordinates for every moment calculation.

The product moment estimate is therefore unchanged:
$\|P_k\|_{L^q(\mu)}\le(qk)^k\max\{U_k,R^{2k-1}\}$.
Each raw marginal of $h$ on at most $k$ coordinates is bounded by
$M(1+\eta)$, so $\|h_S\|_\infty\le2^kM(1+\eta)$.
Every term of $P_k^q$ depends on at most $qk$ coordinates.
Expanding it termwise gives
\[
 \left|\int P_k^q\,d\nu-\int P_k^q\,d\mu\right|
 \le\eta\bigl(2^kN^kM(1+\eta)\bigr)^q\le1.
\]
Orthogonality still gives the exact identity
$U_k^2=\int hP_k\,d\mu=\int FP_k\,d\nu$.
H\"older with respect to $\nu$, the cap on $F$, and
$\int F\,d\nu\le1$ now yields
\[
 U_k^2\le e\left((qk)^k\max\{U_k,R^{2k-1}\}+1\right).
\]
Since $e(qk)^k\le R^{2k}/2$, this inequality rules out $U_k>R^{2k}$:
in that case the maximum is $U_k$, and the right-hand side is strictly
less than $U_k^2$.  This proves the induction, including for the restricted
base measures just described.
\end{proof}

For residues of a uniform integer point, the hypothesis of
Lemma~\ref{sampling:gram} follows directly by counting residues modulo the
product of at most $2b$ of the selected prime powers.  The full residue law
need not be close to product measure.  This distinction prevents an invalid
use of the full product $L^2$ norm of its density.

\subsubsection{Proof of scalar transfer}\label{sampling:scalar-proof}
\begin{proof}[Proof of Proposition~\ref{sampling:scalar-transfer}]
Apply Proposition~\ref{sampling:stability} with its uniform bounds.
Increase polynomial log budgets
when necessary.  We may first work with independent uniform product laws
for $x,V$: domination transfers a sufficiently small failure probability
to any of the allowed sampling laws.  Let $0\le W(x,V)\le1$ be the indicator
of failure.  If the failure probability exceeds $e^{-p}$, then
\begin{equation}\label{sampling:failure-pairing}
 \E_{x,V,t}W(x,V)
 \bigl(fg-(1+\varepsilon)\lambda g\bigr)(x+Vt)>e^{-2p}.
\end{equation}
We will make the upper bound for this expression smaller than $e^{-2p}$.
We first replace the two site functions by their low-degree residue
projections. We then compare the resulting product-coordinate pairing
by conditioning both sides and applying the no-increment hypothesis.

\paragraph{Physical approximation.}
Choose $Q=\exp(\mathcal P(p))$ sufficiently large, at least $e^p$.
For every prime $l\le Q$ use the largest power $q_l\le Q$ as its residue
coordinate.  This collection captures every modulus at most $Q$, and in
particular $D$.  Partition $\mathcal I$ into a fine grid of boxes with
comparable relative side lengths and polynomial log cost, keeping the
boundary of $\mathcal I$ among the grid boundaries.  All site weights and
their approximations are zero outside $\mathcal I$.  Thus stability will
only be applied to restrictions of the original niltest to interior cells.
Extend the grid over a fixed enclosing box for the sampled sites, assigning
zero weights to its exterior cells; its side ratios and log cost have the
initial polynomial bounds.
For $h=g$ or $h=fg$, let $h^\#$ be the density, with respect to the product
reference law of residue coordinates, of the actual $h$-weighted law on
one such cell.  Set, on that cell,
\[
 h^{\mathrm{lo}}(u)=\sum_{|S|\le b}h^\#_S(u_S).
\]
The cutoff $b\ge\log_2Q$ is polynomial and will be chosen after the
conditioning costs below.  The number of terms, each term's cap, and the
sum of caps are $\mathcal B(p)$ once $b$ is fixed.  Products of a polynomial
number of selected moduli have polynomial logarithms.  More explicitly,
let $N$ be the number of selected primes, let
$N_b=\sum_{j\le b}\binom Nj$, and put $C_b=N_b2^b(1+\eta)$.
If a cell has minimum side $J$, every marginal on at most $2b$ prime
coordinates has relative density error at most $O(nQ^{2b}/J)$.
Prescribe this error to be at most $\eta$, with
$\eta C_b^2\le1/10$.  Lemma~\ref{sampling:gram}, followed by the same
termwise comparison for the square, gives
\[
 \E_C(h^{\mathrm{lo}})^2\le2+3\eta C_b^2,
 \qquad \E_C|h-h^{\mathrm{lo}}|^2\le7.
\]
For $m\le Q$ its prime support has size at most $\log_2Q\le b$.
The product projection identity, compared with the actual cell law on
the union of that support and each projection support, therefore gives
\begin{equation}\label{sampling:cell-residue-error}
 \left|\E_C (h-h^{\mathrm{lo}})1_{u\equiv z\pmod m}\right|
 \le\eta C_bm^{-n}.
\end{equation}
Indeed the product projection error against this residue indicator is
exactly zero, and the pointwise relative marginal error contributes its
reference atom mass $m^{-n}$ times $\eta$ and the sum of projection caps.
Only marginals on at most $2b$ primes are used.  The accuracy $\eta$ and
the corresponding side-length cutoff will be chosen after $b$.

We next verify that this replacement is harmless for the averaged sample
in $L^2(x,V)$.  For this estimate alone dominate the uniform spatial laws
by smooth product laws at comparable scales.  In the expansion of the
square, omit the pairs $t,t'$ for which
\[
 \gcd(t'-t)>Q\quad\hbox{or}\quad
 |t'-t|<L/B_0,
\]
where $B_0=\exp(\mathcal P(p))$ is chosen before $b$.
After removing the fixed residue modulo $D$, divisibility counting gives
an upper bound $e^{O_\ell(p)}(Q^{1-\ell}+L^{1-\ell})$ for the first set of pairs,
with an additional negligible diagonal term and an allowed factor from
$D$.  Explicitly, for a nonzero integer difference of size $O(e^pL)$,
the probability that a positive integer $a$ divides all $\ell$ coordinates
is at most $e^{O_\ell(p)}(a^{-1}+L^{-1})^\ell$; summing over
$Q/D<a\le e^{O(p)}L$ gives the assertion.  We later require $L\ge Q$,
which absorbs the finite-length term.  The second set has mass at most
$e^{O_\ell(p)}(B_0^{-1}+L^{-1})^\ell$.
Write $e_h=h-h^{\mathrm{lo}}$.  Summing the cell bounds with their
counting-volume weights gives $\E_{\mathcal I}|e_h|^2\le7$,
independently of the number and sizes of the cells.  Conditional on
$t,t',V$, the translate by $x$ has point mass at most
$e^{\mathcal P(p)}/|\mathcal I|$.  Cauchy--Schwarz therefore bounds its
absolute $e_h(u)e_h(u')$ average by $e^{\mathcal P(p)}$, with a constant
depending only on the initial spatial scales, not on $Q,B_0,b$ or the
grid.  Each discarded set costs this bound times its probability.
Choose $Q,B_0$ to make these contributions negligible.

For a remaining pair put $m=\gcd(t'-t)\le Q$.
The joint location law of $u=x+Vt$ and $u'=x+Vt'$ is zero unless
$u\equiv u'\pmod m$ coordinatewise.  On that congruence its probabilities
are, up to an arbitrarily small normalized error, of the form
\begin{equation}\label{sampling:fiber-density}
 \frac{m^n}{\prod_iH_i^2}\,
 \Phi_{t,t'}(u/H,u'/H),
\end{equation}
where $\Phi_{t,t'}$ has polynomially bounded log cap and Lipschitz constant,
independently of $b$.  Here is a rowwise derivation, retaining the error
needed for the scale bound.  Use smooth compactly supported row weights
$\rho_x(x_i/H_i)\prod_j\rho_v(LV_{ij}/H_i)$, with bounded normalized
derivatives; the allowed scale changes contribute only $e^{\mathcal P(p)}$.
Put $\theta=t/L$, $\Delta=(t'-t)/L$, $s=u_i/H_i$, and $s'=u_i'/H_i$.
Choose $k$ with $|\Delta_k|\ge(\sqrt\ell B_0)^{-1}$.  Eliminating
\[
 z_k=\frac{s'-s-\sum_{j\ne k}\Delta_jz_j}{\Delta_k}
\]
gives the continuous row density
\[
 \Phi_i(s,s')=\frac1{|\Delta_k|}
 \int\rho_x(s-\theta\cdot z)\prod_j\rho_v(z_j)\,dz_{j\ne k}.
\]
Its cap is at most $e^{O_\ell(p)}B_0$ and its Lipschitz constant at most
$e^{O_\ell(p)}(1+B_0)^2$, by differentiation of this formula.

For $d=t'-t$, the integer solutions of $d\cdot V_i=u_i'-u_i$ exist
exactly on the displayed congruence.  Their kernel lattice has covolume
$|d|/m$ in its real hyperplane.  Successive two-entry Bezout column
operations transform $d$ to $(m,0,\ldots,0)$; at most $\ell-1$ operations
with entries bounded by $O_\ell(1+|d|)$ give a kernel basis and a
fundamental region of diameter $O_\ell((1+|d|)^{\ell-1})$.
At the normalized scale $LV_i/H_i$ this diameter is at most
$e^{O_\ell(p)}L^\ell/H_i$.  Partition the hyperplane into these regions.
The error between the smooth weight at a lattice point and its integral
over the region is bounded by its derivative times the diameter; only
regions meeting a fixed bounded enlargement of the support contribute.
After normalization, coarea and the covolume give the row factor
$m/H_i^2$ and error $e^{\mathcal P(p)}L^\ell/H_i$ in its normalized
density.  Taking the product over $n\le p$ rows gives the pointwise law
\begin{equation}\label{sampling:fiber-error}
 1_{u\equiv u'\pmod m}\frac{m^n}{\prod_iH_i^2}
 \left(\Phi_{t,t'}(u/H,u'/H)+O(\epsilon_{\rm fib})\right),
 \qquad
 \epsilon_{\rm fib}\le e^{\mathcal P(p)}\frac{L^\ell}{\min_iH_i}.
\end{equation}
The product changes only polynomial log constants, since each row uses
its own scale $H_i$.  Thus $C'=\ell+1$ suffices; this exponent has no
dependence on $n$, $b$, or the ratios of spatial side lengths.

Choose the first physical grid so that every $\Phi_{t,t'}$ differs by
at most a prescribed $\zeta$ from a constant on each pair of cells.
There is no count loss in using this approximation.  Indeed, CRT gives
$|\{u\in\mathcal I:u\equiv z\pmod m\}|\le2|\mathcal I|m^{-n}$,
so residue-wise Cauchy--Schwarz gives
\[
 \frac{m^n}{|\mathcal I|^2}
 \sum_z\left(\sum_{u\equiv z\pmod m}|e_h(u)|\right)^2
 \le2\E_{\mathcal I}|e_h|^2\le14.
\]
Consequently the grid error and the error in
Equation~\eqref{sampling:fiber-error} cost $O(\zeta+\epsilon_{\rm fib})$.
For the cell-pair constants, Equation~\eqref{sampling:cell-residue-error}
in each cell bounds the sum by
$\|\Phi_{t,t'}\|_\infty(\eta C_b)^2$: the two factors $m^{-n}$ cancel
the fiber factor and the number of residue assignments, and cell volumes
sum as probability weights.  Combining these estimates with the bad-pair
bound yields
\[
 \left\|\E_t e_h(x+Vt)\right\|_2^2
 \le e^{\mathcal P(p)}\left(
 \Pr\{\text{discarded pair}\}+\zeta+\epsilon_{\rm fib}
 +\sup_{\text{retained }t,t'}\|\Phi_{t,t'}\|_\infty(\eta C_b)^2\right).
\]
Thus
\[
 \left\|\E_t(h-h^{\mathrm{lo}})(x+Vt)\right\|_{L^2(x,V)}
 \le\text{prescribed small error}.
\]
Here $Q,B_0$, then $\zeta$ and the first physical grid, are chosen before
$b$.  The fiber-counting and Chinese remainder accuracies can use side
lengths chosen after $b$.

\paragraph{The product-coordinate calculation.}
The physical replacement is now established. Finely mesh the normalized
$x,V,t$ variables so that
$x+Vt$ remains in one physical cell on every retained triple of mesh
boxes.  Triples meeting cell boundaries have small total measure, since
$x$ varies at the full $H_i$ scales; their cost is bounded by the individual
mode caps, which are now available.  On a retained triple, sufficiently
large side lengths allow the $t$ average to be evaluated modulo every
prime support involved, with negligible Chinese remainder error.  Let
$W^\#$ be the product-reference density of the actual $W$-weighted residue
law of $Y=(x,V)$ on its mesh box.  The expressions to compare are
\begin{equation}\label{sampling:bilinear-levels}
 \sum_{|S|\le b}\langle T_Sh^\#_S,W^\#_S\rangle,
 \qquad h=fg,g.
\end{equation}
For a prime coordinate, the reference coupling is
$X_l=x_l+V_lt_l$, with $Y_l=(x_l,V_l)$ uniform and $t_l$ uniform subject
to the imposed $D$-residue.  Both marginals are uniform.  Outside the primes
of $D$, Lemma~\ref{sampling:prime-contraction} gives centered operator norm
at most $l^{-\ell/2}$.  Distinct primes couple independently.

\paragraph{Conditioning the two sides.}
Take $\tau=\exp(-C_0(2+p))$, with $C_0$ sufficiently large.
Choose $r=C_2(2+p)$ next, with $C_2$ large enough to absorb the
$e^{O(2+p)}$ normalized cap losses below using the tail bound $2^{-r}$.
Then choose $C_1$ sufficiently large for the low-level estimates below.
These are choices of fixed constants before any conditioning.
Declare all prime coordinates of $D$ and all $l\le(2+p)^{C_1}$ mandatory.
Their number is polynomially bounded and each has $\log q_l\le\log Q$.
Apply Proposition~\ref{sampling:stability} on the first physical cell,
with these mandatory coordinates, tolerance $\tau$, and further-coordinate
allowance $r+j_*$, where
$j_* = \lceil\log_K(2/\tau)\rceil$ for a fixed $K>1$.
The meaning of mandatory is that every leaf fixes those coordinates.
On a
leaf let $A$ be the fixed coordinates and let $v$ be the conditional
baseline mean.  Also condition on $Y_A$ with the true reference coupling.
The no-increment hypothesis applies on each resulting site cylinder:
its product modulus, physical subbox cost, and hence total log cost are
polynomially bounded.  Indeed these are precisely the subbox and common
modulus slices of Definition~\ref{counting:slice}; the sum of their
logarithmic side losses is polynomial in $p$ since $n\le p$.
Actual conditional means and the corresponding product marginal densities
differ by arbitrarily small relative normalization errors, by Chinese
remainder counting.  Fixing a rare reference atom does not enlarge this
relative error: divide the restricted measure by its reference atom mass,
as in Lemma~\ref{sampling:approx-low-levels}.

Decompose the remaining $W$ density by disjoint cylinder removal.
All means in this removal are exact reference conditional masses of the
weighted density; the underlying physical residual weight stays in $[0,1]$.
If its unconditioned mass is at most $\tau$, retain it as a discard.
Otherwise, among cylinders fixing at most $j_*+r$ coordinates, maximize
their conditional residual mass times $K^{-j}$, where $j$ is the number
fixed.  Each such mass is at most $1+\eta\le2$, whereas the empty
cylinder has mass above $\tau$.  A maximizer therefore has $j\le j_*$.
Any refinement on at most $r$ further coordinates is then included in the
search; maximality implies its conditional mean is at most $K^a$ times
the selected mean, where $a$ coordinates were added.
Remove the entire residual weight on this cylinder and repeat.
A removed cylinder cannot recur, and the number of available cylinders
is $\mathcal B(p)$.  Each retained piece has conditional mean $m\ge\tau$
and, after division by $m$, satisfies the low-level marginal hypothesis
of Lemma~\ref{sampling:approx-low-levels} through level $r$.
The marginal inequalities here are exact, and the normalized physical cap
is at most $1/m$.  Condition also on every possible matching $X$ assignment
on those selected primes.  The already prescribed stability allowance
$j_*+r$ includes these coordinates and $r$ more.  The number of available
cylinders and their inverse reference atom masses are $\mathcal B(p)$:
each uses at most $j_*+r$ prime alphabets of size at most
$Q^{n(\ell+1)}$, from at most $Q$ selected primes.

We justify using the full coordinate coupling on these fixed coordinates
although Equation~\eqref{sampling:bilinear-levels} was truncated.
By linearity split both densities into parts supported on assignments of
a common coordinate set $I$.  A part has the form
$1_{X_I=z}F'(X_{I^c})$ or $1_{Y_I=y}W'(Y_{I^c})$.
For a fixed outside support $S$, summing all inside supports gives exactly
the true coupling weight $\Pr(X_I=z,Y_I=y)$ times the outside pairing.
All inside supports are present when $|S|\le b-|I|$.
The remaining shell, $b-|I|<|S|\le b$, has total absolute cost at most
\begin{equation}\label{sampling:shell-error}
 \mathcal B(p)\,2^{-(b-|I|)}.
\end{equation}
Indeed the centered operator norm on each outside coordinate is at most
$2^{-\ell/2}\le1/2$.  Cauchy--Schwarz over outside supports, with
Lemma~\ref{sampling:gram} for the normalized conditional densities,
bounds their squared-mass sums up to the outside cutoff $b$.
Chinese remainder estimates to order $2b+|I|$ suffice for these uses;
the conditional base measure is divided by its reference atom mass.
The finitely many partial inside sums,
assignment counts, number of pieces, and inverse cylinder probabilities
have polynomial logarithms and are charged to the prefactor.
All those bounds, including $|I|$, were set before $b$.
Choose $b$ last, at least $r+\max|I|$ and sufficiently large to make
Equation~\eqref{sampling:shell-error} negligible.  For discard pieces only
$A$ is conditioned.  No bound on the full product-reference $L^2$ density
has been used.

\paragraph{Completing the comparison.}
On a retained piece, stability bounds every baseline mean obtained
by fixing at most $r$ further prime coordinates by $v+O(\tau)$.
The no-increment hypothesis bounds the corresponding $fg$ means
by $\lambda(v+O(\tau))$, with an additive error at most
$\lambda\tau$.  Thus the natural scales for the two site densities
are $v+\tau$ and $\lambda(v+\tau)$.  Dividing by fixed multiples
of these quantities gives the bounded-marginal hypotheses of
Lemma~\ref{sampling:approx-low-levels}.  The fixed multiples also
absorb the relative CRT normalization errors and make the zeroth
masses at most one.  The normalized logarithmic caps, including
that of $W/m$, are $O(2+p)$: reference conditioning keeps the
unnormalized physical weights bounded by one, while
$m,v+\tau\ge\tau$ and $\lambda\ge e^{-p}$.

For the levels $1,\ldots,r$, the centered contraction
$\kappa=((2+p)^{-C_1})^{\ell/2}$ and
Lemma~\ref{sampling:approx-low-levels} bound the pairing of either
site density with the removed cylinder by
\[
 \text{scale}\cdot m\sum_{k=1}^r
       \kappa^k(C(2+p))^{Ck}.
\]
Choosing $C_1$ after $C_2$ makes this as small a fixed multiple
of $\text{scale}\cdot m$ as required.  For the complete levels
$r<|S|\le b-|I|$, Cauchy--Schwarz and
Lemma~\ref{sampling:gram} give $\kappa^r$ times the product of
the two truncated norms.  Their normalized cap losses are
$e^{O(2+p)}$.  The preceding choice of $C_2$ already controls
this tail using $\kappa\le1/2$, so $r$ need not be increased
after choosing $C_1$.

The zeroth level explains why this comparison suffices.  Its
contribution to the signed pairing is at most
\[
                    -\varepsilon\lambda vm+O(\lambda\tau m),
\]
by \eqref{sampling:no-increment} and stability.  Choose the
contraction errors above to total at most
$\varepsilon\lambda(v+\tau)m/4$ for the two signed terms.
Their sum with the zeroth level is then at most
$O_\varepsilon(\lambda\tau m)$.  This estimate also covers a
vanishing or very small baseline mean; the normalization always
uses $v+\tau$, never $v$ alone.

Discard pieces have cost $O(\sqrt\tau)$ by Cauchy--Schwarz and
Lemma~\ref{sampling:gram}.  Their truncated squared masses are
$O(\tau)$, since their physical caps are bounded and their means
are at most $\tau$, whereas the unnormalized site truncated
squared masses are $O(1)$.  The prescribed Gram errors are chosen
smaller than these bounds.

These errors aggregate without multiplying by the number of cylinders.
Leaf probabilities in the adaptive tree sum to one; $Y_A$ is averaged
conditionally on $X_A$; disjoint removal ensures that the selected
$Y$-cylinder probabilities times their means $m$ sum to a bounded quantity;
and on those coordinates the $X$ assignments use their true conditional
probabilities.  The real mesh triples likewise carry their original
probability weights.  The second, finer mesh only partitions the sample
variables $x,V,t$; the site densities $h^\#$ and their stability trees
still use the original physical $u$ cells, chosen before $b$.  Thus neither
the stability budget nor the no-increment slice cost depends on that finer
mesh.  Its smaller sample boxes require only later Chinese remainder
accuracy and side-length cutoffs.  All cylinder counts are uniform per
mesh triple, and averaging their probability weights incurs no count loss.  Hence, after all approximation errors are included,
\[
 \E W\E_t\bigl(fg-(1+\varepsilon)\lambda g\bigr)(x+Vt)
 \le O(\sqrt\tau)+\text{prescribed approximation errors}.
\]
Choose $C_0$ large and all finite approximation errors sufficiently small
to contradict Equation~\eqref{sampling:failure-pairing}.
This proves Proposition~\ref{sampling:scalar-transfer}.  The choice order is:
$Q,B_0$ and the first physical mesh; stability and cylinder budgets;
$b$; final meshes and accuracies; then sufficiently large $L$ and $H_i$.
Every log cost is polynomial in the original input budget, and the power
$C'=\ell+1$ of $L$ comes only from the $\ell$-dimensional row lattice
calculation; increasing the spatial dimension changes its exponential
prefactor only.  Every later CRT requirement involves a modulus with
polynomial logarithm, so it is met by the final $\mathcal B(p)$ cutoffs
for $L$ and $\min_iH_i$, without enlarging this exponent.
\end{proof}
